%======================================================================
\section{Mathematical and Statistical Foundations}
\label{sec:foundations}

This section states the null models, statistics and decision rules used by
each test, in the form in which they are implemented in the scripts of
Table~\ref{tab:pipeline}. For section $s$ we write $n_s$ for the number of
valid votes (or of votes for the two finalists, when stated), $y_s$ for the
votes of a candidate, $\hat p_s=y_s/n_s$ for the observed share and $\pi_s$
for its expected value under the null hypothesis. $m(s)$ denotes the
municipality and $\ell(s)$ the polling place of the section.

\subsection{Binomial reference and over-dispersion}

The common starting point is that, given the share $\pi_s$ expected for its
electorate, the count of a section is binomial,
\begin{equation}
y_s \mid \pi_s \sim \mathrm{Bin}(n_s,\pi_s),
\label{eq:bin}
\end{equation}
with $\mathrm{E}(y_s)=n_s\pi_s$ and $\mathrm{Var}(y_s)=n_s\pi_s(1-\pi_s)$.
Real sections vary more than (\ref{eq:bin}) allows, because voters are not
assigned to sections at random. Two corrections are used. In the
quasi-binomial approach \cite{wedderburn1974} the variance is multiplied by
a dispersion factor estimated from the Pearson residuals,
\begin{equation}
\hat\varphi=\max\Bigl\{1,\;\frac{1}{N}\sum_{s}\frac{(y_s-n_s\bar\pi_{\ell(s)})^2}{n_s\bar\pi_{\ell(s)}(1-\bar\pi_{\ell(s)})}\Bigr\},
\label{eq:phi}
\end{equation}
where $\bar\pi_\ell$ is the pooled share of the polling place. In the
beta-binomial model \cite{williams1982} $\pi_s$ itself varies around its
mean with intra-section correlation $\rho$, so that
\begin{equation}
\mathrm{Var}(y_s)=n_s\pi_s(1-\pi_s)\bigl[1+(n_s-1)\rho\bigr],
\label{eq:bb}
\end{equation}
and $\rho$ is estimated by the method of moments,
\begin{equation}
\hat\rho=\max\Bigl\{0,\;\frac{\sum_s(r_s^2-1)(n_s-1)}{\sum_s(n_s-1)^2}\Bigr\},
\label{eq:rho}
\end{equation}
with $r_s$ the Pearson residual under (\ref{eq:bin}). Models on the share scale use the
empirical logit $\lambda_s=\ln\{(y_s+\tfrac12)/(n_s-y_s+\tfrac12)\}$, which
is finite at 0\% and 100\% \cite{haldane1956}.

\subsection{Robust scores and multiple testing}

Outliers are measured with the robust $z$-score
\begin{equation}
z_s=\frac{x_s-\operatorname{med}(x)}{1.4826\,\operatorname{med}\bigl|x-\operatorname{med}(x)\bigr|},
\label{eq:zrob}
\end{equation}
where the constant makes the scaled median absolute deviation a consistent
estimator of the standard deviation under normality; median and MAD have
the highest possible breakdown point, so a few extreme sections do not mask
themselves \cite{hampel1974,leys2013}. Thresholds of $|z|\geq4$ or $6$ are
therefore conservative even when $x$ has heavier tails than the normal.

When $m$ hypotheses are tested together, the Benjamini--Hochberg procedure
\cite{benjamini1995} sorts the $p$-values $p_{(1)}\leq\dots\leq p_{(m)}$ and
rejects the $k^\ast$ smallest, with
\begin{equation}
k^\ast=\max\{k: p_{(k)}\leq kq/m\},
\label{eq:bh}
\end{equation}
which keeps the expected proportion of false discoveries below $q$. For a
small family of global tests (eight digit tests, six series of N9) we use
the Bonferroni bound $\alpha/m$ \cite{dunn1961}. All Monte Carlo and
permutation $p$-values are computed as
\begin{equation}
p=\frac{1+\#\{b: T^{(b)}\geq T_{\mathrm{obs}}\}}{B+1},
\label{eq:pmc}
\end{equation}
which is exact under the null and never zero \cite{phipson2010}; its
minimum, $1/(B+1)$, limits the evidence a test with $B$ replicates can
provide.

\subsection{Count regularities (T1--T5, N9)}

\paragraph{T1} With $\pi_s$ taken as the candidate's share in the
\emph{other} sections of the same polling place (or of the municipality, if
the place has a single section), the expected numbers of sections with no
votes and with all votes are
$E_0=\sum_s(1-\pi_s)^{n_s}$ and $E_{100}=\sum_s\pi_s^{n_s}$. Being a sum of
many small, nearly independent probabilities, the observed count is
compared with a Poisson distribution, $p=\Pr\{\mathrm{Pois}(E)\geq O\}$.

\paragraph{T2 and T3} Under Benford's law the second significant digit $d$
has probability \cite{hill1995}
\begin{equation}
\Pr(D_2=d)=\sum_{k=1}^{9}\log_{10}\Bigl(1+\frac{1}{10k+d}\Bigr),\qquad d=0,\dots,9,
\label{eq:benford2}
\end{equation}
and, for counts of two or more digits, the last digit is approximately
uniform. Observed frequencies $O_d$ are compared with expected
frequencies $E_d$ by Pearson's statistic \cite{pearson1900}
\begin{equation}
X^2=\sum_{d=0}^{9}\frac{(O_d-E_d)^2}{E_d}\;\overset{H_0}{\sim}\;\chi^2_9.
\label{eq:chi2}
\end{equation}
Because bounded section sizes distort both references, $E_d$ is also
computed by simulation: each count is redrawn from (\ref{eq:bin}) with
$\pi_s=\hat p_s$, which preserves the size distribution of the sections.

\paragraph{T4 and N9} A count has an integer percentage when
$0<y_s<n_s$ and $100\,y_s\equiv0 \pmod{n_s}$; N9 counts sections whose
share lies within 0.05 percentage points of a multiple of 5\%. In both
cases the statistic $T$ is the number of such sections and its null
distribution is obtained by redrawing every section from
(\ref{eq:bin}) with $\pi_s=\hat p_s$, $B=100$ times in T4 and $B=500$ times
in N9, which is the resampled-kernel idea of Rozenas \cite{rozenas2017}:
spikes produced by small denominators are reproduced, spikes produced by
people are not. The $p$-value follows (\ref{eq:pmc}).

\paragraph{T5} Each section is compared with the rest of its polling place
by the two-proportion statistic
\begin{equation}
z_s=\frac{\hat p_s-\hat p_{-s}}{\sqrt{\bar\pi_\ell(1-\bar\pi_\ell)\bigl(n_s^{-1}+n_{-s}^{-1}\bigr)}},\qquad
\tilde z_s=z_s/\sqrt{\hat\varphi},
\label{eq:t5}
\end{equation}
where $-s$ denotes the other sections of $\ell(s)$ and $\hat\varphi$ is
(\ref{eq:phi}). A section is flagged when $\tilde z_s$ is significant after
(\ref{eq:bh}) with $q=0.001$ and $|\hat p_s-\hat p_{-s}|\geq20$
percentage points.

\subsection{Cross-round and turnout models (T6--T8)}

\paragraph{T6} The runoff share of a finalist is regressed on the
first-round composition of the same section, $\mathbf c_s=(c_{s1},\dots,c_{sK})$,
the shares of every candidate and of blank and null votes (one category
omitted, since they sum to one), with municipality fixed effects removed by
the weighted within transformation (a double dot denotes the deviation from
the municipal mean weighted by $n_s$, the regression weight):
\begin{equation}
\ddot p^{(2)}_s=\sum_{k=1}^{K-1}\beta_k\,\ddot c_{sk}+e_s,
\label{eq:t6}
\end{equation}
The $R^2$ of (\ref{eq:t6}) is the
within-municipality $R^2$ reported in Section~\ref{sec:results}; the scaled residuals
$e_s\sqrt{n_s}$ are converted to robust scores (\ref{eq:zrob}). Turnout
change is $\Delta C_s=\ln(V_{2,s}/V_{1,s})$, centred on the municipality.

\paragraph{T7} The fingerprint is the two-dimensional histogram of turnout
$t_s=V_s/A_s$ ($A_s$ eligible voters) and leader share. In the model of
Klimek et al.\ \cite{klimek2012}, ballot stuffing moves a fraction of the
abstentions and of the opponent's votes to the leader in a subset of units,
which produces a second mode near $(t,\hat p)=(1,1)$ and a positive
association between $t_s$ and $\hat p_s$ within regions; both are checked.

\paragraph{T8} Goodman's ecological regression \cite{goodman1953}
estimates, by state and without intercept, how first-round votes transfer to
each finalist:
\begin{equation}
V^{A}_{2,s}=\sum_{j\in\{A,B,O,\mathit{BN}\}}\gamma^{A}_{j}V_{1,s}^{j}+u_s,
\label{eq:goodman}
\end{equation}
and analogously for $B$, where $O$ are the other candidates and
$\mathit{BN}$ blank and null votes. The predicted share is
$\tilde p_{2,s}=\hat V^A_{2,s}/(\hat V^A_{2,s}+\hat V^B_{2,s})$ and the
unexplained shift is $\Delta P_s=100(p_{2,s}-\tilde p_{2,s})$. The
transfer coefficients $\gamma$ are valid proportions only if transfer rates
are constant within the state; values outside $[0,1]$ signal that this
assumption fails. The association between $\Delta C$ and $\Delta P$ is
measured by Spearman's rank correlation \cite{spearman1904},
$\rho_S=1-6\sum_i d_i^2/\{N(N^2-1)\}$ with $d_i$ the rank differences, and by
decile means with normal 95\% intervals. Among the $M$ sections with extreme
$\Delta C>0$ and $\Delta P$, the number $X$ that shifted toward the first
finalist is tested against $X\sim\mathrm{Bin}(M,\tfrac12)$, the distribution
expected if neither side benefited.

\subsection{Fixed effects and clustered errors (N4--N8, N13)}

N4--N7 estimate linear models of the form $y_s=\mathbf x_s^\top\boldsymbol\beta+\alpha_{g(s)}+\varepsilon_s$
with weights $w_s$ and fixed effects $\alpha_g$ for the polling place (N6)
or the municipality (N4, N5, N7; also N13). The weighted within transformation
$\ddot y_s=y_s-\sum_{g(i)=g(s)}w_iy_i/\sum_{g(i)=g(s)}w_i$ (and likewise for
$\mathbf x_s$) removes $\alpha_g$, and
$\hat{\boldsymbol\beta}=(\ddot X^\top W\ddot X)^{-1}\ddot X^\top W\ddot{\mathbf y}$.
Because sections of the same municipality share unobserved shocks, the
variance is estimated with the cluster-robust sandwich \cite{cameron2015}
\begin{equation}
\widehat{\mathrm V}(\hat{\boldsymbol\beta})=c\,\mathbf H^{-1}
\Bigl(\sum_{h=1}^{G}\mathbf g_h\mathbf g_h^\top\Bigr)\mathbf H^{-1},
\label{eq:sandwich}
\end{equation}
with $\mathbf H=\ddot X^\top W\ddot X$ and cluster scores
$\mathbf g_h=\sum_{s\in h}w_s\ddot{\mathbf x}_s\hat u_s$ summed over the
$G$ municipalities, and small-sample factor
$c=G(n-1)/\{(G-1)(n-k)\}$. Two-sided $p$-values use the normal
approximation and the BH correction (\ref{eq:bh}).

In N7 the model is fitted on the empirical logit $\lambda_s$; the fitted
probability is $\hat\pi_s=\operatorname{logit}^{-1}(\mathbf x_s^\top\hat{\boldsymbol\beta}+\hat\alpha_{m(s)})$,
over-dispersion is estimated from (\ref{eq:bb}), and each section is scored
by $z_s=(y_s-n_s\hat\pi_s)/\{n_s\hat\pi_s(1-\hat\pi_s)[1+(n_s-1)\hat\rho]\}^{1/2}$.
The model is checked by simulation \cite{rubin1984,gelman1996}: $R=200$
replicated data sets $y^{(r)}_s\sim\text{Beta-Bin}(n_s,\hat\pi_s,\hat\rho)$
give the reference distribution of discrepancy statistics $T(\mathbf y)$
(sections at 0\%, above 90\% and below 10\%, and the within-municipality
standard deviation), with two-sided predictive $p$-value
$p=\min\{1,2\min(\Pr[T(\mathbf y^{(r)})\geq T(\mathbf y)],\Pr[T(\mathbf y^{(r)})\leq T(\mathbf y)])\}$.
In N8 the change $\delta_s=\lambda^{2026}_s-\lambda^{2022}_s$ of the same
section is centred on the municipal median and scored with (\ref{eq:zrob})
by state.

\subsection{Spatial dependence (N10)}

Let $z_i$ be a standardized residual averaged over polling place $i$ and
$w_{ij}=1/k$ if $j$ is one of the $k=8$ nearest polling places of $i$ in
the same state (zero otherwise). Global and local Moran statistics are
\cite{moran1950,anselin1995}
\begin{equation}
I=\frac{\sum_i z_i\sum_j w_{ij}z_j}{\sum_i z_i^2},\qquad
I_i=z_i\sum_j w_{ij}z_j .
\label{eq:moran}
\end{equation}
$I>0$ means that neighbouring places have similar residuals. The local
statistic is assessed by conditional permutation: the neighbours of $i$
are replaced $B=199$ times by places drawn at random, and $p_i$ follows
(\ref{eq:pmc}), with (\ref{eq:bh}) at $q=0.01$. Since $p_i\geq1/(B+1)=0.005$
and (\ref{eq:bh}) over $m\approx9\times10^4$ places requires
$p_{(k)}\leq kq/m$, no place can be significant unless at least
$0.005\,m/q\approx4.6\times10^4$ places reach that minimum together; this
is why the null result of N10 is
uninformative (Section~\ref{sec:extended}).

\subsection{Classification and multimodality (N11, N12)}

\paragraph{N11} A random forest \cite{breiman2001,wright2017} estimates,
from synthetic sections, the probability $\hat f(\mathbf x)$ that a section
with features $\mathbf x$ was manipulated. Discrimination is measured by the area under the ROC curve, which
equals the Mann--Whitney probability that a manipulated section receives a
higher score than a clean one \cite{mann1947,hanley1982}:
\begin{equation}
\mathrm{AUC}=\frac{R_1-n_1(n_1+1)/2}{n_1n_0},
\label{eq:auc}
\end{equation}
with $R_1$ the rank sum of the $n_1$ manipulated sections among all
$n_1+n_0$ scores. The decision threshold is the $(1-\alpha)$ quantile of
the scores of held-out clean sections, $\tau=\hat F_0^{-1}(1-\alpha)$ with
$\alpha=0.001$, so that among $N$ clean real sections about $\alpha N$
would be flagged by chance.

\paragraph{N12} Hartigan's dip \cite{hartigan1985} is the distance between
the empirical distribution function $F_n$ of the section shares of a
municipality and the closest unimodal distribution,
\begin{equation}
D(F_n)=\inf_{G\in\mathcal U}\;\sup_x\bigl|F_n(x)-G(x)\bigr|,
\label{eq:dip}
\end{equation}
where $\mathcal U$ is the class of unimodal distribution functions. Its
$p$-value is taken from the null distribution under the uniform, the least
favourable unimodal case, and the municipalities are corrected with
(\ref{eq:bh}) at $q=0.05$.

\subsection{Integrity of the machine files (N1, N2)}

N1 is an identity test: for every section and office, the vector of votes
recounted from the RDV must equal the vector printed in the BU, and the BU
must equal the published CSV. N2 relies on two primitives. SHA-512
\cite{fips180} maps a file to a 512-bit digest $h=H(F)$; a change in any bit
of $F$ changes $h$ unpredictably, so equality of the recomputed digest with
the signed one shows that the file is the one that was signed. The signature
on the digest list is verified with the public key embedded in the
machine certificate. For ECDSA over P-521 \cite{fips186,johnson2001}, with
base point $G$ of order $q$, public key $Q$, message digest $e$ and signature
$(r,s)$, the verifier computes
\begin{equation}
w=s^{-1}\bmod q,\quad (x_1,y_1)=(ew\bmod q)\,G+(rw\bmod q)\,Q,
\label{eq:ecdsa}
\end{equation}
and accepts iff $x_1\equiv r\pmod q$. For EdDSA \cite{bernstein2012,rfc8032}
on the Edwards curve E-521, $x^2+y^2=1-376014\,x^2y^2$ over
$\mathbb F_{2^{521}-1}$ \cite{aranha2013}, a signature $(R,S)$ on message
$M$ under public key $A$ is accepted iff
\begin{equation}
[S]B=R+[k]A,\qquad k=H(R\,\|\,A\,\|\,M),
\label{eq:eddsa}
\end{equation}
with $B$ the base point. Forging either signature without the private key
requires solving the elliptic-curve discrete logarithm problem, so a valid
signature binds the files to the key of the certificate they carry; it does
not, by itself, show that the key belongs to a genuine machine.
